\section*{Bab \ref{ch:b1:comparative-planetology} --- Planetologi komparatif}

\begin{solution}{exo:b1:comparative-planetology:1}
$V = \tfrac43\pi(\qty{6.9911e7}{m})^3 = \qty{1.4313e24}{m^3}$ dan
$\bar\rho = \num{1898.13e24}/\num{1.4313e24} = \qty{1326}{kg.m^{-3}}$, hanya
1.33 kali massa jenis air, meskipun Jupiter 318 kali lebih massif daripada Bumi: ia tersusun sebagian besar dari unsur-unsur paling ringan.
\end{solution}

\begin{solution}{exo:b1:comparative-planetology:2}
Merkurius (jari-jari \qty{2440}{km}, \qty{5429}{kg.m^{-3}}) dan Venus
(\qty{6052}{km}, \qty{5243}{kg.m^{-3}}) kecil dan lebih padat daripada batuan: \omterm{def:b1:comparative-planetology:planet-types}{planet terestrial}, dari batuan dan besi. Saturnus (\qty{60268}{km},
\qty{687}{kg.m^{-3}}) sangat besar dan lebih ringan daripada air: \omterm{def:b1:comparative-planetology:planet-types}{raksasa gas}, sebagian besar hidrogen dan helium. Uranus (\qty{25559}{km}, \qty{1270}{kg.m^{-3}}) adalah raksasa tetapi jauh lebih kecil daripada Saturnus dan hampir dua kali lebih padat: \omterm{def:b1:comparative-planetology:planet-types}{raksasa es}, yang air, metana, dan amonianya lebih berat daripada hidrogen.
\end{solution}

\begin{solution}{exo:b1:comparative-planetology:3}
$T_\mathrm{eq} = (2601.3\times0.23/(4\times\num{5.6704e-8}))^{1/4} =
(\num{2.638e9})^{1/4} = \qty{226.6}{K}$. Venus menerima 1.9 kali lebih banyak sinar Matahari daripada Bumi, tetapi awannya memantulkan \qty{77}{\percent} darinya: ia hanya menyerap $2601.3\times0.23 = \qty{598}{W.m^{-2}}$, dibandingkan dengan $1361\times0.706 = \qty{961}{W.m^{-2}}$ untuk Bumi.
\end{solution}

\begin{solution}{exo:b1:comparative-planetology:4}
Dari yang paling seragam sampai yang paling terkonsentrasi: Bulan, Mars, Bumi, Jupiter, Saturnus. Massa titik di pusat memiliki $I = 0$, sehingga $C = 0$.
\end{solution}

\begin{solution}{exo:b1:comparative-planetology:5}
$T = 15.945421\times86400 = \qty{1377684}{s}$;
$GM = 4\pi^2(\num{1.22187e9})^3/\num{1377684}^2 = \qty{3.794e16}{m^3.s^{-2}}$;
$M = \num{3.794e16}/\num{6.6743e-11} = \qty{5.685e26}{kg}$, dalam
\qty{0.03}{\percent} dari \qty{568.32e24}{kg}.
\end{solution}

\begin{solution}{exo:b1:comparative-planetology:6}
$S = 1361\times(149.6/228.0)^2 = \qty{585.9}{W.m^{-2}}$ di Mars dan
$1361\times(149.6/778.5)^2 = \qty{50.26}{W.m^{-2}}$ di Jupiter, sesuai
dengan tabel (selisih kecil di Mars berasal dari pembulatan jarak rata-rata).
\end{solution}

\begin{solution}{exo:b1:comparative-planetology:7}
$S \propto d^{-2}$ dan $T_\mathrm{eq} \propto S^{1/4}$, sehingga $T_\mathrm{eq}
\propto d^{-1/2}$. Untuk Jupiter, $T_\mathrm{eq} = (50.26\times0.657/(4\times
\num{5.6704e-8}))^{1/4} = \qty{109.8}{K}$. Planet yang memancarkan lebih banyak daripada yang diserapnya pasti memiliki sumber panas internal: Jupiter melepaskan panas yang tersimpan sejak pembentukannya, atau yang dihasilkan ketika ia perlahan mengerut.
\end{solution}

\begin{solution}{exo:b1:comparative-planetology:8}
Venus: $737 - 226.6 = \qty{510}{K}$, \qty{225}{\percent} dari $T_\mathrm{eq}$.
Bumi: $T_\mathrm{eq} = (1361\times0.706/(4\sigma))^{1/4} = \qty{255.1}{K}$,
$288 - 255.1 = \qty{33}{K}$, \qty{13}{\percent}. Mars: $214 - 209.8 =
\qty{4}{K}$, \qty{2}{\percent}.
\end{solution}

\begin{solution}{exo:b1:comparative-planetology:9}
Dengan $\mathrm dm = 4\pi\rho r^2\,\mathrm dr$: $I = \int_0^R\tfrac23r^2\,4\pi
\rho r^2\,\mathrm dr = \tfrac{8\pi}{15}\rho R^5$, dan $MR^2 = \tfrac43\pi\rho
R^5$, sehingga $C = \tfrac{8}{15}\cdot\tfrac34 = \tfrac25$. Kulit bola berongga memiliki seluruh massanya pada $r = R$: $I = \tfrac23MR^2$, $C = \tfrac23$.
\end{solution}

\begin{solution}{exo:b1:comparative-planetology:10}
$\tfrac52C\bar\rho = 2.5\times0.35\times5429 = 4750.4$;
$x^2 = (4750.4 - 3300)/(5429 - 3300) = 1450.4/2129 = 0.6812$, $x = 0.8254$,
jari-jari inti $0.8254\times2439.7 = \qty{2014}{km}$. $x^3 = 0.5623$ dan
$\rho_c = 3300 + 2129/0.5623 = \qty{7086}{kg.m^{-3}}$. Inti memuat
$0.5623\times7086/5429 = \qty{73}{\percent}$ massanya. Inti Merkurius
mencapai \qty{83}{\percent} dari jari-jarinya, dibandingkan dengan $3485/6371 =
\qty{55}{\percent}$ untuk Bumi: Merkurius sebagian besar inti, dengan mantel batuan yang tipis.
\end{solution}

\begin{solution}{exo:b1:comparative-planetology:11}
(a) Fungsi $t\mapsto t^4$ adalah cembung, sehingga
$\bigl(\tfrac{T_1 + T_2}{2}\bigr)^4 \le \tfrac{T_1^4 + T_2^4}{2} =
T_\mathrm{eq}^4$, dengan kesamaan hanya bila $T_1 = T_2$. (b) $T_1^4 =
2\times270^4 - 100^4 = \num{1.0529e10}$, $T_1 = \qty{320}{K}$, dan suhu
rata-ratanya $(320 + 100)/2 = \qty{210}{K}$, \qty{60}{K} di bawah
$T_\mathrm{eq}$.
\end{solution}

\begin{solution}{exo:b1:comparative-planetology:12}
$m = 4\pi R^2p/g = 4\pi\times(\num{6.0518e6})^2\times\num{9.2e6}/8.9 =
\num{4.602e14}\times\num{1.034e6} = \qty{4.76e20}{kg}$: 93 kali massa atmosfer Bumi, dan \num{9.8e-5} massa Venus. Bumi menyimpan karbon dalam jumlah yang sebanding, tidak di udaranya tetapi dalam batuan karbonat (batu kapur) dan di samudra, tempat air cair memungkinkannya tersimpan.
\end{solution}

\begin{solution}{pb:b1:comparative-planetology:1}
\textbf{1.} $T = 0.31891\times86400 = \qty{27553.8}{s}$.

\textbf{2.} $GM = 4\pi^2\times(\num{9.378e6})^3/\num{27553.8}^2 =
\qty{4.289e13}{m^3.s^{-2}}$.

\textbf{3.} $M = \num{4.289e13}/\num{6.6743e-11} = \qty{6.426e23}{kg}$,
dalam \qty{0.14}{\percent} dari \qty{6.4169e23}{kg} dalam tabel.

\textbf{4.} $\num{1.06e16}/\num{6.426e23} = \num{1.6e-8}$: Phobos dapat diabaikan.

\textbf{5.} $v = 2\pi a/T = 2\pi\times\num{9.378e6}/27553.8 =
\qty{2138}{m.s^{-1}}$, sekitar \qty{2.1}{km.s^{-1}}.

\textbf{6.} $T = \qty{7.65}{h}$, 3.2 kali lebih pendek daripada rotasi
Mars: Phobos mengitari lebih cepat daripada tanah berputar di bawahnya, sehingga ia mendahului rotasi. Dilihat dari Mars, ia terbit di barat dan terbenam di timur.

\textbf{7.} $V = \tfrac43\pi(\num{3.3895e6})^3 = \qty{1.6312e20}{m^3}$.

\textbf{8.} $\bar\rho = \num{6.426e23}/\num{1.6312e20} = 3939$, yaitu
\qty{3940}{kg.m^{-3}}.

\textbf{9.} $3940/5514 = 0.715$. Bumi, yang lebih besar dan lebih massif, memampatkan bagian dalamnya sendiri lebih kuat; dan inti besinya merupakan bagian yang lebih besar dari planet itu.

\textbf{10.} $g = GM/R^2 = \num{4.289e13}/(\num{3.3895e6})^2 =
\qty{3.73}{m.s^{-2}}$, \qty{38}{\percent} dari gravitasi Bumi.

\textbf{11.} $v_\mathrm{esc} = \sqrt{2GM/R} = \sqrt{2\times\num{4.289e13}/
\num{3.3895e6}} = \qty{5.03}{km.s^{-1}}$.

\textbf{12.} $T_\mathrm{eq} = (586.2\times0.75/(4\times\num{5.6704e-8}))^{1/4}
= \qty{209.8}{K}$.

\textbf{13.} $214 - 209.8 = \qty{4}{K}$: pemanasan rumah kaca yang lemah.

\textbf{14.} $m = p/g = 636/3.733 = \qty{170}{kg.m^{-2}}$; untuk Bumi
$101400/9.82 = \qty{10326}{kg.m^{-2}}$, 61 kali lebih banyak.

\textbf{15.} $4\pi R^2\times170.4 = \num{1.444e14}\times170.4 =
\qty{2.46e16}{kg}$.

\textbf{16.} Di atas setiap meter persegi terdapat gas enam puluh kali lebih sedikit daripada di Bumi, sehingga serapan inframerah, meskipun oleh karbon dioksida, lemah; dan udara Mars hanya menyimpan sangat sedikit uap air, penyerap utama lainnya.

\textbf{17.} Mars yang seragam akan memiliki $C = 0.4$; nilai terukur 0.366 menunjukkan bahwa massa jenisnya bertambah ke arah pusat: Mars memiliki inti yang padat.

\textbf{18.} $\bar\rho = \rho_m + (\rho_c - \rho_m)x^3$.

\textbf{19.} $C = \tfrac25[\rho_m + (\rho_c - \rho_m)x^5]/\bar\rho$.

\textbf{20.} Dari 18 dan 19, $(\rho_c - \rho_m)x^3 = \bar\rho - \rho_m$ dan
$(\rho_c - \rho_m)x^5 = \tfrac52C\bar\rho - \rho_m$; membagi keduanya menghasilkan $x^2$.

\textbf{21.} $\tfrac52C\bar\rho = 2.5\times0.366\times3940 = 3605.1$;
$x^2 = (3605.1 - 3500)/(3940 - 3500) = 105.1/440 = 0.2389$, $x = 0.4887$.
$x^3 = 0.1167$ dan $\rho_c = 3500 + 440/0.1167 = \qty{7269}{kg.m^{-3}}$.

\textbf{22.} $x^3\rho_c/\bar\rho = 0.1167\times7269/3940 = 0.215$, yaitu
\qty{21.5}{\percent} massanya.

\textbf{23.} Dengan $\rho_m = 3450$: $x^2 = 155.1/490 = 0.3165$, $x = 0.5626$,
jari-jari inti \qty{1907}{km}, $\rho_c = \qty{6202}{kg.m^{-3}}$. Dengan $\rho_m =
3550$: $x^2 = 55.1/390 = 0.1413$, $x = 0.3759$, \qty{1274}{km},
$\rho_c = \qty{10894}{kg.m^{-3}}$. Perubahan \qty{1.4}{\percent} pada massa jenis mantel yang diasumsikan menggeser jari-jari inti sebanyak beberapa ratus kilometer: massa jenis rata-rata dan $C$ saja tidak dapat menentukan inti Mars dengan tepat, itulah sebabnya seismometer dikirim ke sana.

\textbf{24.} $0.4887\times3389.5 = \qty{1657}{km}$: jari-jari inti Mars dari model adalah \textbf{sekitar \qty[text-series-to-math=true]{1660}{km}}, dalam rentang seismik terbaru \qty{1650}{km} sampai \qty{1675}{km} dan jauh di bawah \qty{1830}{km} dari tabel planet --- kesesuaian yang ditunjukkan soal 23 sebagian merupakan keberuntungan dalam pemilihan $\rho_m$.
\end{solution}
